\scanpage{33}
la prop. 3, \deleted{\unsure{et 2}} montrant le lien exact entre $\infty$-polygônes réguliers et repr. affines du groupe diédral infini, enfin la prop 4 (où dans b), \inserted{et c),} il ne faut plus parler de questions de divisibilité, puisque $n$ est ici remplacé par $+\infty$\ldots)

\uline{NB} Ici encore, pour une classe d'isom. de réal. géom. \inserted{régulières} \deleted{\unsure{de|des}} polygône combinatoire, \unsure{déterminé|donné}, \inserted{\unsure{et les assoc.,}} \deleted{épinglée \unsure{(biff.)} il y a \unsure{bien}} un ``module'' $\alpha\in k$, -- maintenant $\alpha\in k$ n'est plus soumis à aucune condition, cyclotomique ou autre --) il y a lieu d'écrire ``\uline{la}'' \unsure{espa.|repr.} affine $E$ correspondant \uline{au} représentant épinglé (aut. unique) Plus \unsure{gén$^{\rm t}$|général} et plus joliment ou intrinsèque, si $\Pi$ est un poly combinatoire fini ou infini, \unsure{notant|notons} que la cat. de ses réalisations géom. régulières est \uline{rigide} (ses objets, n'ont d'auto. autres que l'identité) il y a lieu pour \unsure{la} classe d'iso (déterminée tjrs par le $=\alpha$) d'écrire $E\inserted{=E(\Pi)}$ -- \deleted{étant entendu} que ce sera une
% slanted
\margin{C'est un des intérêts techniques d'éliminer \deleted{\unsure{au départ}} au début un \unsure{choix} de la théorie\ldots}
% archive sheet number: 32
